\chapter{Verificación}
\label{ch:ch5_verification}

% [BORRADOR v1] Material de partida: rtl/tb/*, rtl/run_tb.py, analysis/ctr_drbg_ref.py,
% analysis/verifica_top.py, firmware/verificador/, docs/mapa_registros.md.
% Sección de puesta en marcha en la FPGA pendiente del hardware.

\section{Estrategia: vectores, simulación, verificador independiente}
\label{sec:verif_strategy}

Un motor criptográfico que solo se verifica contra sí mismo no demuestra
nada: un error en la S-box se reproduciría idéntico en el modelo usado
para comprobarlo. Por eso, la verificación de este trabajo se apoya en
tres pilares que no comparten implementación.

\begin{enumerate}
\item \textbf{Vectores oficiales.} Cada primitiva se comprueba frente a
los vectores de prueba publicados con su norma: FIPS~197 y SP~800-38A
para AES, RFC~4493 y el fichero de ejemplos del NIST para CMAC, y los
ficheros de respuesta del CAVP para CTR\_DRBG (con y sin función de
derivación, con y sin resembrado). Son pruebas de respuesta conocida y
zanjan la corrección del algoritmo.
\item \textbf{Modelos de referencia independientes.} Se escribió en
Python un modelo exacto al bit del generador, con su propio AES,
derivado de la norma sin mirar el VHDL. Primero se valida frente a los
vectores del CAVP, y solo entonces se usa para generar los valores
esperados para la disposición fija de entradas del hardware y para
recalcular las etiquetas del banco de pruebas de sistema. Dos
implementaciones que coinciden valen más que una que pasa las pruebas.
\item \textbf{Un verificador independiente en el lazo.} En hardware, el
ESP32 recalcula con mbedtls, y con su propio acelerador AES, cada
resultado que produce el motor, de modo que la comprobación final la
realiza una tercera implementación, sin relación con las otras, sobre un
bus real.
\end{enumerate}

Dos reglas de trabajo se aplican en todo momento, ambas heredadas del
trabajo experimental previo del autor: la salida de una compilación,
simulación o medida nunca se filtra antes de haberla leído entera --- un
filtro oculta justo el aviso que importaba --- y los fallos se registran
completos, nunca se descartan. Dos veces en este trabajo un filtro sobre
el log del simulador ocultó un resultado, y las dos veces la solución fue
quitar el filtro.

\section{Modelos de referencia en Python}
\label{sec:refmodels}

\texttt{ctr\_drbg\_ref.py} implementa el cifrado AES-128, la función de
derivación basada en cifrado por bloques, CBC-MAC y CTR\_DRBG con y sin
función de derivación, en Python puro y a partir del texto de
SP~800-90A. El AES comprueba con aserciones los dos vectores de FIPS~197
al importarse; el DRBG se prueba frente a los ficheros de respuesta del
CAVP siguiendo el procedimiento CAVS (instanciar, resembrado opcional,
dos llamadas de generación, comparar la segunda): \textbf{960 casos
AES-128, todos superados}. Se encontró un defecto en el analizador de
ficheros y no en el modelo --- una cabecera de sección debe cerrar el
caso de prueba pendiente antes de que cambien los parámetros; si no, el
último caso de cada bloque de parámetros se evalúa con los parámetros del
bloque siguiente --- y merece la pena dejarlo anotado porque su síntoma,
un único fallo en count~14 de cada bloque, parecía exactamente un error
del algoritmo.

Un segundo script, \texttt{gen\_drbg\_vectores.py}, hace pasar al modelo
validado por la secuencia que admite el hardware (instanciar con 32
bytes de entropía y 16 de nonce, generar dos veces, resembrar, generar)
y emite un paquete VHDL con los valores esperados; si alguna vez cambia
la disposición de las entradas del hardware, ese script es el único
sitio donde hay que cambiarla. Un tercero, \texttt{verifica\_top.py},
recalcula AES-CMAC a partir del mismo AES y se valida a sí mismo frente
a RFC~4493 antes de comprobar nada.

\section{Bancos de pruebas y resultados}
\label{sec:testbenches}

Todo el VHDL se simula con GHDL~6 mediante un único script de ejecución
que analiza las fuentes en orden de dependencias y ejecuta cada banco,
imprimiendo la salida del simulador sin filtrar. La
Tabla~\ref{tab:tb} resume los siete bancos y su estado en el momento de
la redacción.

\begin{table}[htbp]
\centering
\caption{Bancos de pruebas, qué comprueban y su resultado (GHDL 6.0, VHDL-2008).}
\label{tab:tb}
\small
\begin{tabular}{@{}lp{9.2cm}c@{}}
\toprule
Banco & Comprobaciones & Resultado \\
\midrule
Oscilador en anillo & arranque/parada, $f = 1/(2N t_{pd})$, suelo nulo, $\sigma_T = \sqrt{2N}\,\sigma_{\mathrm{stage}}$ & 4/4 \\
Esclavo I2C & direccionamiento, dirección ajena, escritura multibyte, lectura con inicio repetido & 14/14 \\
Núcleo AES-128 & FIPS~197 B, C.1; SP~800-38A F.1.1 (3 bloques); cadena de 1000 bloques & 6/6 \\
AES-CMAC & RFC~4493 ejemplos 1--4; repetibilidad & 5/5 \\
CTR\_DRBG & instanciación, 2$\times$generación, resembrado, contador, reinstanciación, rechazo sin semilla & 13/13 \\
Tests de salud & fuente sana, atascada, sesgada al 70\,\%, atascada en el complemento, retención/borrado & 10/10 \\
Motor completo & secuencia del ESP32 por I2C: id, semilla (divisor por debajo del mínimo rechazado, divisor congelado mientras está ocupado), clave, registro aleatorio libre de la clave, 2 retos, 2.ª clave, captura, clave oculta & 18/18 \\
\bottomrule
\end{tabular}
\end{table}

Parte de lo que encontraron los bancos merece contarse, porque un
capítulo de verificación que solo informa de aciertos no ha verificado
gran cosa.

\paragraph{Modelo del anillo.} Las primitivas LUT usadas para simular el
anillo se modelaron al principio con una regla estricta de propagación
de valores desconocidos: cualquier entrada desconocida daba una salida
desconocida. Al comienzo de una simulación todas las señales son
desconocidas, así que el anillo nunca salía de ese estado y nunca
arrancaba. Una NAND real con su entrada de habilitación a cero da un uno
a la salida sea cual sea la otra entrada; el modelo ahora evalúa todas
las combinaciones compatibles con las entradas conocidas y resuelve la
salida cuando todas coinciden. Sin esto, todo el banco de pruebas de
sistema habría sido imposible.

\paragraph{Esclavo I2C.} Dos errores de diseño reales: el puntero de
registro no avanzaba en las escrituras multibyte, de modo que los bytes
consecutivos sobrescribían la misma dirección; y en las lecturas el bit
más significativo se presentaba en el bus un flanco de reloj demasiado
tarde, lo que desplazaba todos los bytes. El primero se corrigió
incrementando el puntero un estado después del pulso de escritura, de
forma que la dirección sigue siendo válida durante el pulso; el segundo,
cargando el registro de desplazamiento y presentando su primer bit en el
mismo flanco de bajada que libera el bit de reconocimiento. Un tercer
defecto estaba en el propio banco: inicializar la memoria de registros
desde el proceso de estímulos creaba un segundo driver sobre la misma
señal, y todo bit en conflicto se resolvía a desconocido.

\paragraph{Tests de salud.} La fuente ``sana'' del primer banco, un
registro de desplazamiento inicializado con un solo uno, empezaba con
una racha de 27 ceros y una primera ventana al 76\,\%; el detector saltó.
El detector tenía razón.

\paragraph{Fuga de la clave por el registro aleatorio.} El hallazgo más
grave no salió de un banco, sino de releer el diseño antes de
entregarlo. La orden que generaba la clave escribía la salida completa
de 256 bits del DRBG en el registro de salida aleatoria, que se puede
leer sin el pin de prueba; los primeros 128 de esos bits \emph{eran} la
clave. El banco de sistema incluso imponía esa igualdad mediante una
aserción, como ``comprobación''. La corrección separa ambas cosas: la
generación de la clave no deja nada legible, y una orden nueva llena el
registro aleatorio a partir de una generación independiente, después de
que el DRBG haya actualizado su estado. Ahora el banco comprueba con una
aserción lo contrario, que el registro no contiene la clave en ninguna
de sus mitades, y la comprobación cruzada en Python hace lo mismo.

\paragraph{Motor completo.} Con los anillos a su velocidad física, el
banco de sistema, que genera casi un milisegundo de tráfico I2C, no
terminaba en veinte minutos de CPU: cada transición del anillo cuesta
dos sorteos aleatorios en el modelo de jitter. Solo en ese banco el
retardo de etapa se fija diez veces mayor, con el mismo jitter relativo,
y así se indica en el propio banco. Las dieciocho comprobaciones se
superan, y la clave, los dos retos y sus etiquetas se escriben en un
fichero que comprueba \texttt{verifica\_top.py}: las dos etiquetas son
iguales al CMAC recalculado por el AES independiente en Python, el
registro de salida aleatoria legible no contiene la clave en ninguna de
sus mitades, y dos generaciones dan claves distintas. Este es el primer
punto en que una clave nacida del ruido de osciladores en anillo,
acondicionada, generada y usada para autenticar, queda confirmada por
una implementación que no comparte nada con el VHDL.
Dos de las dieciocho comprobaciones se añadieron cuando dos obligaciones
que el mapa de registros imponía al anfitrión se trasladaron al
hardware: una orden de siembra emitida con el divisor por debajo de su
mínimo ahora se rechaza, y una escritura en el registro del divisor
mientras se ejecuta una orden se ignora. Una regla que solo impone la
documentación es una regla que un anfitrión descuidado incumple; el
banco ahora prueba ambas. Además, hubo que indicarle al banco que se
detuviera: con los anillos funcionando tras el reset, una simulación
cuyos estímulos han terminado nunca se queda sin eventos, y la primera
repetición consumió cuarenta y cinco minutos de CPU después de haber
superado su última comprobación.

\section{Firmware verificador del ESP32}
\label{sec:firmware}

El firmware es una aplicación ESP-IDF~5 para el ESP32 clásico que se
comunica con el motor a través del mapa de registros del
Apéndice~\ref{app:regs} y ofrece una consola por la UART. Sus órdenes se
corresponden con el plan de verificación: \texttt{id} lee la
identificación y el estado; \texttt{sembrar} y \texttt{generar} lanzan
las órdenes largas y consultan el indicador de ocupado; en modo de
prueba, \texttt{generar} además relee la clave (aprovisionamiento) y
comprueba que es igual a los dieciséis primeros bytes del bloque
generado; \texttt{reto N} realiza $N$ rondas de reto-respuesta con
nonces de \texttt{esp\_random()}, recalculando cada etiqueta con
\texttt{mbedtls\_cipher\_cmac} y contando coincidencias y discrepancias,
y las discrepancias se imprimen completas; \texttt{crudo N} vuelca $N$
fragmentos de bits en bruto y \texttt{esp N} vuelca bloques del propio
generador del ESP32, ambos en el formato en tramas protegido con CRC32
que usó el autor en su trabajo anterior, para que la misma cadena de
procesado estadístico trate las dos fuentes. CMAC no viene habilitado
por defecto en la compilación de mbedtls de ESP-IDF y se activa en la
configuración por defecto del proyecto; la consola funciona a
115\,200~baudios, la única velocidad que resultó fiable con el puente
USB de la placa en campañas anteriores, y el nivel de log se fija en
solo errores para que las líneas de log no puedan corromper las tramas.
El firmware compila con ESP-IDF~5.3.2 en una imagen de 255~kB;
ejecutarlo queda a la espera del hardware.

\section{Puesta en marcha en la FPGA}
\label{sec:bringup}

% [PENDIENTE DE HARDWARE] Proyecto de Vivado, restricciones (DONT_TOUCH, pblocks,
% ALLOW_COMBINATORIAL_LOOPS, set_disable_timing), primer bitstream: ¿oscilan los
% anillos y a qué frecuencia?; después, el ESP32 en el lazo.
Pendiente del hardware. Las restricciones que necesitan los anillos se
conocen por la bibliografía y por la documentación del fabricante ---
\texttt{DONT\_TOUCH} en las celdas y redes de cada anillo (a diferencia
de \texttt{KEEP}, sobrevive a la implementación),
\texttt{ALLOW\_COMBINATORIAL\_LOOPS} en una red de cada lazo,
\texttt{set\_disable\_timing} para romper el lazo de cara al análisis
estático de tiempos, y un bloque de emplazamiento por anillo para que
anillos idénticos ni se fusionen ni se coloquen uno junto a otro --- y
el primer bitstream tiene un único propósito: leer los contadores de
frecuencia y confirmar que todos los anillos oscilan. Antes de eso no se
afirma nada sobre la entropía.
